% AI-written, AI cross-reviewed; no human sign-off; not a claim that RH is proved
% Atlas unit W14. Local manuscript v3, item 5.50, heading on printed PDF p. 53.
% UpperTheory.V.Odd.thmBcFree_of
% paper2/split_p013.tex:286–312. Verbatim excerpt; SHA-256 991ee82ada4257fe2ab2acd397908f35221389a43d84cf0ad869658c58dbe4e2.
% Original TeX uses the macros and cross-reference keys of the source manuscript.

\begin{formalstatement}[From special values to the minus part]\label{sub:W14}
Assume the following character estimate. For every abelian CM field
$F$, compatible algebraic closure, and cyclotomic action of
$(\mathbb Z/N\mathbb Z)^\times$ on $F$ ($N>0$), take an odd prime
$\ell\nmid[F:\mathbb Q]$, a finite DVR
$O\subset\overline{\mathbb Q}_\ell$ over $\mathbb Z_\ell$, $k\ge1$,
and a character $\psi$ of $\operatorname{Gal}(F/\mathbb Q)$ valued in
$O^\times$. Write $\chi$ for the induced character modulo $N$.
If $\chi(-1)=(-1)^{k+1}$, the $O$-length of the $\psi$-eigenspace
of the localization kernel with twist $-k$ is at most
$\operatorname{ord}_O(a)$ for every $a\in O$ representing the
denominator-cleared primitive Bernoulli value
$\operatorname{den}(N,k+1)B_{\rm prim}(k+1,\chi)$.
Then there exists a real $C$ such that for every abelian datum $M$,
$e,t\ge2$, prime pair $q,d$ as in W12 and CM field $E$ of type
$(M,e,q,d)$, all its odd-prime minus lengths $h_\ell^-(E,k)$ for
$1\le k<t$ are finite. Put $n=[E:\mathbb Q]$, $f=\operatorname{cond}E$,
$o=\deg M$. For every finite set $T$ of odd primes,
\[
\begin{split}
 \sum_{\ell\in T}\log\ell\sum_{1\le k<t}h_\ell^-(E,k)
 &\le 3nt^2\log(tf)\\
 &\quad+C\!\sum_{\substack{p\mid o\\p\ {\rm odd\ prime}}}
                  pn^2t^2\log^2(tpf).
\end{split}
\]
\end{formalstatement}\subnode{W14}
